\chapter{Nguyên lý ATPG}
\label{sec:p1-nguyen-ly-atpg}
\textbf{Mô hình lỗi và điều kiện phát hiện.}
Lỗi kẹt mức (stuck-at) buộc một đường tín hiệu luôn bằng 0 (SA0) hoặc 1
(SA1); giả thiết lỗi đơn xét từng lỗi riêng. Lỗi trên thân dây (stem) tác
động mọi nhánh đi ra, còn lỗi nhánh (branch) chỉ tác động một ngõ vào cổng.
Mô hình này đơn giản, phù hợp ATPG mức cổng nhưng không mô tả mọi khuyết tật:
lỗi cầu nối (bridging) liên quan các dây nối ngoài ý muốn, còn lỗi trễ
(delay) liên quan đáp ứng không kịp thời gian~\cite[Ch.~1]{p1_wang2006}.
Để phát hiện lỗi phải \emph{kích hoạt} giá trị đối nghịch tại vị trí lỗi,
\emph{lan truyền} khác biệt đến PO và \emph{justify} các yêu cầu bằng gán PI
nhất quán. Logic năm giá trị dùng $0,1,X,D,D'$, với
$D=(1,0)$, $D'=(0,1)$ theo thứ tự (mạch tốt, mạch lỗi); $X$ là chưa biết.
Trên c17, theo thứ tự PI $(1,2,3,6,7)$, mẫu $(X,1,0,X,X)$ cho lỗi
\texttt{11/SA0} tạo $11=D$, $16=D'$, $10=1$, $22=D$.
Mọi cách điền ba giá trị $X$ đều phát hiện lỗi này.

\textbf{Danh sách và rút gọn lỗi.}
Nếu chỉ xét hai lỗi trên mỗi net, c17 có $2(5+6)=22$ lỗi; thêm lỗi nhánh
sẽ thay đổi mẫu số. Gọi $T(f)$ là tập vector phát hiện lỗi $f$.
Hai lỗi tương đương về phát hiện khi $T(f_1)=T(f_2)$; tương đương đáp ứng
mạnh hơn, yêu cầu cùng đáp ứng lỗi ở mọi PI. Với NAND $z=\overline{ab}$,
$a/\mathrm{SA0}$, $b/\mathrm{SA0}$ và $z/\mathrm{SA1}$ tương đương tại cổng.
Trên c17, net 1 chỉ nối đến cổng 10 nên có thể giữ một đại diện cho
$1/\mathrm{SA0}$ và $10/\mathrm{SA1}$: cả hai buộc $10=1$, do đó có cùng
đáp ứng PO trên cả 32 vector. Không áp dụng quy tắc cục bộ này trực tiếp
cho stem có nhiều nhánh. Với dominance, điều kiện
$T(f_a)\subseteq T(f_b)$ cho phép giữ $f_a$: mọi mẫu phát hiện $f_a$ cũng
phát hiện $f_b$. Ví dụ NAND riêng có $T(a/\mathrm{SA1})=\{01\}$ nằm trong
$T(z/\mathrm{SA0})=\{00,01,10\}$. Cần ghi chiều bao hàm để tránh nhầm
lỗi đại diện~\cite[pp.~12--14]{p1_wang2006}.

\textbf{Mô phỏng và độ bao phủ.}
Mô phỏng lỗi (fault simulation) kiểm tra một bộ mẫu trên mạch tốt và mạch lỗi;
kiểu serial xét từng lỗi, còn kiểu parallel có thể mã hóa nhiều mạch lỗi vào
các bit của một từ máy. Fault dropping loại lỗi đã được phát hiện khỏi
danh sách còn phải xét. Với danh sách lỗi gốc $F$, độ bao phủ là
$FC=|F_{\mathrm{det}}|/|F|\times100\%$; phải đếm mỗi lỗi một lần và ghi rõ
stem/branch, trước/sau collapsing. Nếu báo cáo test coverage loại các lỗi
đã chứng minh không thể phát hiện, mẫu số là $|F|-|F_{\mathrm{untestable}}|$.
Không tìm được mẫu trong giới hạn tài nguyên là \texttt{ABORTED}, không phải
\texttt{UNTESTABLE}. Trong mạch tổ hợp, lỗi không có vector phát hiện được
gọi là untestable; lỗi dư thừa là một trường hợp liên quan logic dư thừa.
Theo quy ước nhóm, điền $X=0$ trước mô phỏng nhị phân và kiểm tra lại mẫu.

\textbf{Khả năng điều khiển, quan sát và SCOAP.}
SCOAP ước lượng chi phí đặt net về 0/1 ($CC0,CC1$) và quan sát nó ($CO$);
số nhỏ hơn biểu thị dễ hơn, không phải xác suất. Đặt $CC0=CC1=1$ tại PI,
$CO=0$ tại PO. Với $z=\operatorname{NAND}(a,b)$~\cite[pp.~41--42]{p1_wang2006}:
\[
\begin{aligned}
CC0(z)&=CC1(a)+CC1(b)+1, & CC1(z)&=\min(CC0(a),CC0(b))+1,\\
CO(a\!\to\!z)&=CO(z)+CC1(b)+1, & CO(s)&=\min_{\text{nhánh }j}CO(s_j).
\end{aligned}
\]
Tính $CC$ từ PI đến PO rồi $CO$ theo chiều ngược lại thu được
Bảng~\ref{tab:p1-scoap}. Chẳng hạn $CC0(11)=1+1+1=3$,
$CC1(11)=2$; $CO(11)=\min(3+1+1,3+1+1)=5$.
Do fanout hội tụ lại (reconvergence), SCOAP bỏ qua tương quan giữa các nhánh;
đây là heuristic định hướng tìm kiếm, không chứng minh một lỗi có thể phát hiện.
\begin{table}[htbp]
\centering\small
\caption{SCOAP c17, tính trên net/stem theo quy tắc trong giáo trình.}
\label{tab:p1-scoap}
\begin{tabular}{@{}lrrrrrrrrrrr@{}}
\toprule
Net & 1 & 2 & 3 & 6 & 7 & 10 & 11 & 16 & 19 & 22 & 23\\
\midrule
$CC0$ & 1 & 1 & 1 & 1 & 1 & 3 & 3 & 4 & 4 & 5 & 5\\
$CC1$ & 1 & 1 & 1 & 1 & 1 & 2 & 2 & 2 & 2 & 4 & 5\\
$CO$  & 5 & 6 & 5 & 7 & 6 & 3 & 5 & 3 & 3 & 0 & 0\\
\bottomrule
\end{tabular}
\end{table}

\textbf{Phương pháp sinh mẫu.}
Sinh mẫu ngẫu nhiên ít tốn công tìm kiếm nhưng khó xử lý lỗi khó phát hiện.
ATPG tất định hướng tìm kiếm vào từng lỗi: D-algorithm (Roth, 1966) thao tác
trên các đường tín hiệu; PODEM (Goel, 1981) quyết định tại PI; FAN
(Fujiwara và Shimono, 1983) khai thác cấu trúc fanout để giảm tìm kiếm
\cite[Ch.~4]{p1_wang2006}. Nhóm chọn PODEM để trình bày objective, backtrace,
imply và backtrack; chưa mặc định đây là thuật toán nhanh nhất.
Hình~\ref{fig:p1-atpg-flow} mô tả luồng xử lý chung.
\begin{figure}[htbp]
\centering
% Sơ đồ do P1 vẽ; các gói TikZ được khai báo ở main.tex.
\begin{tikzpicture}[>=Stealth,node distance=4mm,
  every node/.style={font=\footnotesize,align=center},
  box/.style={draw,rounded corners,text width=3.5cm,minimum height=8mm}]
\node[box] (select) {Chọn lỗi chưa xử lý};
\node[box,right=8mm of select] (search) {Kích hoạt, lan truyền,\\justify; backtrack khi xung đột};
\node[box,right=8mm of search] (result) {Phân loại kết quả tìm kiếm};
\node[box,below=of result] (stop) {Hết tìm kiếm: UNTESTABLE\\Hết giới hạn: ABORTED};
\node[box,below=of search] (sim) {Có mẫu: mô phỏng kiểm chứng;\\lưu mẫu và fault dropping};
\draw[->] (select)--(search);
\draw[->] (search)--(result);
\draw[->] (result)--(stop);
\draw[->] (result.south west)--(sim.north east);
\draw[->] (sim.west)-|(select.south);
\end{tikzpicture}

\caption{Luồng ATPG tất định và kiểm chứng mẫu.}
\label{fig:p1-atpg-flow}
\end{figure}
